\documentclass[10pt,twocolumn]{article}
\usepackage[T1]{fontenc}
\usepackage[utf8]{inputenc}
\usepackage{lmodern}
\usepackage[margin=1.8cm,top=2.4cm,bottom=2cm,columnsep=0.6cm]{geometry}
\usepackage{fancyhdr}
\usepackage{titlesec}
\usepackage{booktabs}
\usepackage{amsmath,amssymb,amsfonts}
\usepackage{microtype}
\usepackage{xcolor}
\usepackage{mdframed}
\usepackage{parskip}
\usepackage{enumitem}
\usepackage{hyperref}

\definecolor{rulered}{RGB}{180,30,30}
\definecolor{boxgray}{RGB}{245,245,245}
\definecolor{keywordgray}{RGB}{100,100,100}

\pagestyle{fancy}
\fancyhf{}
\fancyhead[L]{\textsc{\small Claude Mag}}
\fancyhead[C]{\small\textit{Machine Learning Research Digest}}
\fancyhead[R]{\small 2026-11-09}
\fancyfoot[C]{\small\thepage}
\renewcommand{\headrulewidth}{0.8pt}

\titleformat{\section}{\normalsize\bfseries\color{rulered}}{}{0pt}{}%
  [\vspace{-4pt}\color{rulered}\rule{\columnwidth}{0.4pt}]
\titlespacing{\section}{0pt}{8pt}{4pt}

\hypersetup{colorlinks=true,urlcolor=rulered,linkcolor=black}

\begin{document}
\setlength{\parindent}{0pt}
\setlength{\parskip}{4pt}

{\small\color{keywordgray}\textbf{Keywords:}
  reinforcement learning \textbullet{}
  reward machines \textbullet{}
  vision-language models \textbullet{}
  non-Markovian tasks}\par\vspace{4pt}

{\LARGE\bfseries\raggedright
  RM-VLM: From Natural-Language Task Descriptions to
  Vision-Grounded Reward Machines}\par\vspace{4pt}

{\large\itshape\raggedright
  An LLM Builds the Finite Automaton; a VLM Grounds Its Propositions;
  the Agent Earns Credit for History-Dependent Task Progress}\par\vspace{6pt}

{\small\textbf{By} Roko Parac, Christos Ziakas, Leo Ardon}\par\vspace{2pt}

{\small\color{keywordgray}arXiv:2611.08308 \textbullet{}
  cs.AI \textbullet{} submitted 6 Nov 2026}
\par\vspace{2pt}\noindent{\color{rulered}\rule{\columnwidth}{0.6pt}}\vspace{6pt}

Most VLM-based reward methods assume that reward depends only on the
current observation. RM-VLM handles tasks that require history-dependent
credit by having an LLM construct a reward machine from a text
description and a VLM ground each proposition in image observations.
Evaluated on grid navigation, robotic manipulation, and 3D object
collection, RM-VLM consistently outperforms existing VLM-reward methods
on task success.

\begin{mdframed}[backgroundcolor=boxgray,linecolor=rulered,linewidth=1pt,
    innerleftmargin=6pt,innerrightmargin=6pt,
    innertopmargin=6pt,innerbottommargin=6pt]
\textbf{\small AT A GLANCE}\par\vspace{4pt}
\begin{description}[leftmargin=0pt,labelsep=4pt,itemsep=2pt,parsep=0pt,
    font=\small\bfseries]
  \item[Problem:] VLM reward methods assume Markovian tasks; real tasks
    often require ordered sub-goals whose credit depends on history.
  \item[Why it matters:] Without history-aware reward, RL agents
    fail to distinguish correct-order completions from incorrect-order
    ones and cannot learn ordered, multi-step tasks reliably.
  \item[Method:] LLM decomposes task description into propositions and
    builds a reward machine (finite automaton); VLM grounds propositions
    in image observations; VLM value estimator provides dense shaping.
  \item[Result:] Consistently higher task success than all baselines
    across three task domains.
  \item[Evidence:] Grid-world navigation, robotic manipulation sequence,
    3D object collection; baselines: VLM-RM, CLIP-reward, TAMER+VLM.
\end{description}
\end{mdframed}

\section{The Big Picture}

Reinforcement learning agents historically required hand-crafted reward
functions that encode both the goal and the intermediate milestones.
VLM-based reward models automate this by scoring image observations
against a natural-language goal description: an agent that moves closer
to the goal receives a higher reward from the VLM, without any manual
reward engineering.

This works well for Markovian goals where the reward depends only on
the current state. It breaks for \emph{ordered} tasks: ``pick up the
red block, then place it in the green zone, then open the door.'' A VLM
reward that evaluates only the current image cannot distinguish whether
the red block was placed \emph{before} or \emph{after} the green zone
interaction. An agent that learns to shortcut to the door without
completing earlier sub-goals will receive high reward on the final
state.

Reward machines (Icarte et al., 2018) are the standard formalism for
non-Markovian task structure: a finite automaton tracks which sub-goals
have been achieved and transitions on logical propositions evaluated in
the environment. RM-VLM combines the LLM's ability to decompose
natural-language task descriptions with the VLM's ability to ground
propositions in images.

\section{Why Existing Approaches Fall Short}

\textbf{VLM scalar reward} (VLM-RM) scores the current image against
the full task description. It cannot distinguish sub-goal ordering and
provides identical reward to agents that complete sub-goals in any order.

\textbf{CLIP-reward} embeds both image and text goal in CLIP space and
uses cosine similarity as reward. It has the same Markovian limitation
and additionally fails to track incremental progress toward structured
sub-goals.

\textbf{TAMER with VLM feedback} uses intermittent human-style feedback
modeled by a VLM. The feedback is episodic and noisy; it does not
provide dense, automaton-guided credit for individual transitions.

\textbf{Handcrafted reward machines} (prior RM work) require a domain
expert to specify the propositions and automaton structure manually ---
the bottleneck RM-VLM eliminates.

\section{Core Mechanism}

\textbf{Step 1: Reward machine construction.}
The task description $\mathcal{T}$ is passed to an LLM with a
structured prompt requesting:
\begin{itemize}[itemsep=2pt,parsep=0pt]
  \item A set of Boolean propositions $\mathcal{P} = \{p_1, \ldots,
    p_n\}$ over environment entities.
  \item A finite automaton
    $M = (Q, q_0, q_f, \delta, \rho)$ where
    $\delta: Q \times 2^\mathcal{P} \to Q$ is the transition function
    and $\rho: Q \times Q \to \mathbb{R}$ assigns reward to transitions.
\end{itemize}
The LLM outputs $M$ as a JSON object, which is validated and compiled.

\textbf{Step 2: VLM grounding.}
At each step, the VLM receives the current image observation $o_t$ and
each proposition $p_i$, returning $P(p_i{=}\text{true} \mid o_t)$.
The resulting Boolean vector advances $M$'s state through $\delta$,
triggering the reward $\rho(q_t, q_{t+1})$ on any transition.

\textbf{Step 3: Dense reward shaping.}
To alleviate sparse rewards in long-horizon tasks, a VLM-based value
estimator $V_{\text{VLM}}$ scores the image against the proposition
set of each pending transition, producing a shaped bonus:
\[
  r^+_t = \lambda \sum_{p_i \in \text{pending}(q_t)}
    P(p_i{=}\text{true} \mid o_t),
\]
where $\lambda$ is a shaping coefficient and $\text{pending}(q_t)$
is the set of propositions required to advance from $q_t$.

\section{Technical Formulation}

The reward machine $M = (Q, q_0, q_f, \delta, \rho)$ defines a
non-Markovian MDP $\mathcal{M} = (S, A, T, R_M)$ where the total
reward at step $t$ is:
\[
  R_M(s_t, a_t) = \rho(q_t, q_{t+1}) + r^+_t,
\]
with $q_{t+1} = \delta(q_t, L(o_t))$ and
$L(o_t) \in 2^\mathcal{P}$ the labelling function implemented by the
VLM. The augmented state $(s_t, q_t)$ is Markovian in
$\mathcal{M}$, so standard policy-gradient methods (PPO, SAC) apply.

The labelling function uses a threshold $\tau = 0.5$:
\[
  L(o_t) = \{p_i : P(p_i{=}\text{true} \mid o_t) \ge \tau\}.
\]

\section{Learning or Inference Procedure}

\begin{enumerate}[itemsep=2pt,parsep=0pt]
  \item Submit task description $\mathcal{T}$ to LLM to obtain $M$
    as JSON.
  \item Validate $M$ (check that $q_f$ is reachable, no duplicate
    states, all propositions well-formed).
  \item Initialize policy $\pi_\theta$ with PPO; use
    $(s_t, q_t)$ as the augmented state.
  \item At each environment step, query VLM for $L(o_t)$, advance $M$,
    and compute $R_M(s_t, a_t)$.
  \item Train until convergence (5M environment steps per domain).
\end{enumerate}

No fine-tuning of the LLM or VLM is performed; both are used frozen.

\section{What the Guarantee Says}

No formal guarantee is claimed. The empirical claim is consistent
task-success superiority over all baselines in all three domains
tested. No sample complexity bound is given; the 5M-step budget is
task-specific and may not be sufficient for harder domains.

\section{Experimental Findings}

\begin{itemize}[itemsep=2pt,parsep=0pt]
  \item \textbf{Grid-world navigation:} RM-VLM task success 87\%
    vs.\ VLM-RM 61\%, CLIP-reward 54\%, TAMER+VLM 70\%.
  \item \textbf{Robotic manipulation sequence:} RM-VLM 74\% vs.\
    VLM-RM 38\%, CLIP 29\%, TAMER+VLM 51\%; largest gap because
    out-of-order completions are highest in this domain.
  \item \textbf{3D object collection:} RM-VLM 81\% vs.\ VLM-RM 58\%.
  \item \textbf{Dense shaping:} Removing the shaped bonus
    $r^+_t$ reduces robotic manipulation success from 74\% to 52\%,
    confirming its role in long-horizon credit propagation.
\end{itemize}

\section{Ablations and Interpretation}

\textbf{LLM construction errors:} On 15\% of randomly sampled task
descriptions, the LLM produces a machine with at least one missing
transition or unreachable state. These cases are caught by validation
and trigger a retry; the retry resolves 90\% of errors.

\textbf{VLM grounding accuracy:} Replacing the VLM labelling function
with ground-truth oracle labels improves manipulation success from 74\%
to 88\%, establishing an upper bound from grounding accuracy limitations.

\textbf{Threshold $\tau$:} Grid search over $\tau \in \{0.3, 0.5, 0.7\}$
shows $\tau = 0.5$ optimal across all three domains; $\tau = 0.3$
causes false-positive transitions; $\tau = 0.7$ misses valid ones.

\textbf{Shaping coefficient $\lambda$:} $\lambda = 0.1$ is optimal;
larger values cause the agent to maximize the shaped bonus rather than
completing transitions, a known reward-shaping pathology.

\section*{Reference}

Roko Parac, Christos Ziakas, Leo Ardon.
\textbf{RM-VLM: From Natural-Language Task Descriptions to
Vision-Grounded Reward Machines.}
arXiv:2611.08308.
\url{https://arxiv.org/abs/2611.08308}

\end{document}
